\begin{lythuyet}
\loidan{Cùng Cherry xem vì sao một chữ số có thể mang giá trị khác nhau khi đổi chỗ.}
\quytac{Hệ thập phân}{Dùng mười chữ số $0;\,1;\,2;\,\dots;\,9$ để ghi số. Cứ $10$ đơn vị ở một hàng thì bằng $1$ đơn vị ở hàng liền trước nó. Giá trị của mỗi chữ số phụ thuộc vào vị trí của nó trong số.}
\congthuc{Cấu tạo thập phân của số}{$\overline{ab}=a\cdot 10+b$ ($a\ne 0$) \\ $\overline{abc}=a\cdot 100+b\cdot 10+c$ ($a\ne 0$)}
\minhhoa{gia-tri-theo-hang}{2485}
\luuy{Phân biệt chữ số hàng trăm và số trăm: số $2\,485$ có chữ số hàng trăm là $4$, còn số trăm là $24$.}
\vidu{Trong số $6\,052$, chữ số hàng nghìn và chữ số hàng chục đóng góp bao nhiêu?}{$6$ ở hàng nghìn có giá trị $6\,000$. $5$ ở hàng chục có giá trị $50$.}
\thungay{Trong $3\,672$, giá trị của chữ số $6$ là bao nhiêu?}{$600$}{$36$}{$6$ ở hàng trăm nên có giá trị $6\cdot 100=600$.}{$36$ là số trăm trong $3\,672$. Riêng chữ số $6$ ở hàng trăm có giá trị $600$.}

\loidan{Đồng hồ có thể dùng chữ thay cho chữ số. Em thử ghép các mảnh để viết số La Mã nhé!}
\quytac{Số La Mã}{$\mathrm I=1$, $\mathrm V=5$, $\mathrm X=10$. Hai chữ đặc biệt: $\mathrm{IV}=4$, $\mathrm{IX}=9$. \\ Chữ có giá trị nhỏ hơn đứng sau thì cộng thêm: $\mathrm{XII}=10+1+1=12$. Các chữ $\mathrm I$, $\mathrm X$ không đứng liền nhau quá ba lần.}
\minhhoa{ghep-so-la-ma}{24}
\vidu{Viết số $27$ bằng số La Mã.}{$27=10+10+5+1+1$ nên $27=\mathrm{XXVII}$.}
\thungay{Số La Mã $\mathrm{XVIII}$ có giá trị là bao nhiêu?}{$18$}{$13$}{Sau $\mathrm{X}$ là $\mathrm{V}$ và ba chữ $\mathrm{I}$, nên ta cộng $10+5+1+1+1=18$.}{Khi cộng $10$ với ba chữ $\mathrm{I}$, em còn cần cộng $5$ của chữ $\mathrm{V}$.}
\end{lythuyet}
